\documentclass[11pt,a4paper]{article}
\usepackage[margin=2.5cm]{geometry}
\usepackage{amsmath,amssymb,amsthm,bm}
\usepackage{booktabs}
\usepackage{array}
\usepackage{algorithm}
\usepackage{algpseudocode}
\usepackage{hyperref}

\newtheorem{proposition}{Proposition}
\newtheorem{remark}{Remark}

\newcommand{\Fb}{\bar{\mathbf F}}
\newcommand{\Pb}{\bar{\mathbf P}}
\newcommand{\sigb}{\bar{\bm\sigma}}
\newcommand{\Pt}{\mathbf P^{\mathrm t}}
\newcommand{\sigt}{\bm\sigma^{\mathrm t}}
\newcommand{\F}{\mathbf F}
\newcommand{\C}{\mathbf C}
\newcommand{\Pk}{\mathbf P}
\newcommand{\I}{\mathbf I}
\newcommand{\x}{\mathbf x}
\newcommand{\X}{\mathbf X}
\newcommand{\w}{\mathbf w}
\newcommand{\f}{\mathbf f}
\newcommand{\U}{\mathbf U}
\newcommand{\T}{\mathbf T}
\newcommand{\G}{\mathbf G}
\newcommand{\m}{\mathbf m}
\newcommand{\tr}{\operatorname{tr}}

\title{Stress and traction control in a Lagrange-reduced lattice finite-element model:\\
nominal and Cauchy stress loading of periodic cells, and dead-load boundary tractions}
\date{}

\begin{document}
\maketitle

\begin{abstract}
We describe how mechanical loads are prescribed in force (rather than displacement) in a
two-dimensional lattice finite-element model in which the strain energy of every triangle is
evaluated after Lagrange (\(\mathrm{GL}(2,\mathbb Z)\)) reduction of its metric. Two settings are
treated. For periodic cells, the macroscopic deformation gradient of the cell becomes an unknown
conjugate to a prescribed macroscopic stress. A prescribed nominal (first Piola--Kirchhoff) stress
is a dead load and leads to the minimization of an enthalpy in the joint space of internal
coordinates and cell deformation. A prescribed true (Cauchy) stress is a follower load that admits
no potential; it is enforced exactly by a fixed-point iteration on the nominal target, each pass
being an enthalpy minimization. For finite specimens, dead-load nominal tractions on boundary
segments and rigid load-controlled grips are introduced through external work potentials. For both
settings we give the equilibrium equations, the work-conjugate strain measures used to detect
plastic strain bursts, the treatment of rigid-body modes and of the gauge rotation, and the
detection of the limit load past which no equilibrium exists. Finally, Newton-type alternatives
based on the Cauchy--Born, relaxed (Schur-complement) and full Hessian tangents are formulated and
compared with the first-order method.
\end{abstract}

%=====================================================================
\section{Discrete energy}
\label{sec:energy}
%=====================================================================

The reference configuration is a triangulation of the lattice sites \(\{\X_a\}\). Each triangle
\(e\) with vertices \(a\in e\) carries constant shape-function gradients \(\nabla N_a^e\), so that
its deformation gradient is
\begin{equation}
  \F_e=\sum_{a\in e}\x_a\otimes\nabla N_a^e ,\qquad \C_e=\F_e^{\mathsf T}\F_e .
  \label{eq:Fe}
\end{equation}
The strain-energy density per unit reference area is obtained from a lattice (pair or many-body)
energy \(\Phi\) evaluated on the reduced metric,
\begin{equation}
  W(\C)=\frac{1}{A_{\mathrm{cell}}}\,\Phi\!\left(\m^{\mathsf T}\C\,\m\right),\qquad
  \m=\m(\C)\in\mathrm{GL}(2,\mathbb Z),
  \label{eq:W}
\end{equation}
where \(\m(\C)\) is the Lagrange reduction of \(\C\) and \(A_{\mathrm{cell}}\) is the area of the
lattice unit cell, which makes \(W\) an energy per unit area. By construction \(W\) is invariant
under \(\C\mapsto \mathbf n^{\mathsf T}\C\mathbf n\) for all \(\mathbf n\in\mathrm{GL}(2,\mathbb Z)\),
i.e.\ under lattice-invariant shears, which is what allows dislocations to nucleate and glide
without any reconnection of the underlying lattice. The first Piola--Kirchhoff stress of a
triangle is
\begin{equation}
  \Pk_e=\frac{\partial W}{\partial\F}(\F_e)
  =\frac{2}{A_{\mathrm{cell}}}\,\F_e\,\m\,
  \frac{\partial\Phi}{\partial\C}\big(\m^{\mathsf T}\C_e\m\big)\,\m^{\mathsf T},
  \label{eq:Pe}
\end{equation}
and the total elastic energy is
\begin{equation}
  E_{\mathrm{el}}=\sum_e A_e\,W(\C_e),\qquad
  \frac{\partial E_{\mathrm{el}}}{\partial\x_a}=\sum_{e\ni a}A_e\,\Pk_e\,\nabla N_a^e .
  \label{eq:E}
\end{equation}
Here \(A_e\) is the reference area of triangle \(e\). Because \(\nabla N^e_a\) is fixed by the
reference element, \(A_e\) must be the area of that reference element,
\(A_e=\tfrac12\,|\det[\nabla N^e_2,\nabla N^e_3]|^{-1}\), and not the area of the triangle formed
by the reference positions of the current vertices: after a mesh reconnection the two differ, and
using the latter corrupts the energy and the stress of the reconnected triangles.

%=====================================================================
\section{Periodic cells: kinematics and macroscopic stress}
\label{sec:pbc}
%=====================================================================

\subsection{Cell deformation and internal coordinates}

A periodic cell has reference area \(A_0\) and is deformed by a homogeneous macroscopic
deformation gradient \(\Fb\). Atomic positions are written as
\begin{equation}
  \x_a=\Fb\,\w_a ,
  \label{eq:xw}
\end{equation}
where \(\w_a\) are internal (cell) coordinates. A triangle that straddles the cell boundary uses
periodic images of its vertices, \(\x_a+\Fb\,\T_a^e\), with \(\T^e_a\) a fixed translation by a
reference period vector. Substituting in \eqref{eq:Fe},
\begin{equation}
  \F_e=\Fb\,\G_e,\qquad \G_e=\sum_{a\in e}\big(\w_a+\T^e_a\big)\otimes\nabla N^e_a ,
  \label{eq:FeG}
\end{equation}
where \(\G_e\) does not depend on \(\Fb\). The energy is therefore a function
\(E_{\mathrm{el}}(\w,\Fb)\), and its partial derivatives at fixed \(\w\) and at fixed \(\Fb\) are
\begin{align}
  \frac{\partial E_{\mathrm{el}}}{\partial\w_a}
    &=\Fb^{\mathsf T}\,\frac{\partial E_{\mathrm{el}}}{\partial\x_a},
  \label{eq:dEdw}\\
  \frac{\partial E_{\mathrm{el}}}{\partial\Fb}
    &=\sum_e A_e\,\Pk_e\,\G_e^{\mathsf T}
     =\Big(\sum_e A_e\,\Pk_e\,\F_e^{\mathsf T}\Big)\Fb^{-\mathsf T}.
  \label{eq:dEdF}
\end{align}
The matrix \(\mathbf V=\sum_e A_e\Pk_e\F_e^{\mathsf T}\) is the (Kirchhoff) virial of the cell. It
is accumulated during the same element loop that computes the energy and the nodal forces, so the
derivative with respect to the cell costs nothing extra.

\subsection{Macroscopic stresses}

The macroscopic nominal and Cauchy stresses are defined by work conjugacy with the cell
deformation,
\begin{equation}
  \Pb=\frac{1}{A_0}\,\frac{\partial E_{\mathrm{el}}}{\partial\Fb}
     =\frac{1}{A_0}\,\mathbf V\,\Fb^{-\mathsf T},\qquad
  \sigb=\frac{1}{\bar J}\,\Pb\,\Fb^{\mathsf T}=\frac{\mathbf V}{A_0\,\bar J},
  \qquad \bar J=\det\Fb .
  \label{eq:macrostress}
\end{equation}
Since \(\Pk_e\F_e^{\mathsf T}=J_e\bm\sigma_e\) and \(A_eJ_e=a_e\) is the current area of the
triangle, \(\sigb=\sum_e a_e\bm\sigma_e/(A_0\bar J)\) is the current-area average of the element
Cauchy stresses, provided \(A_0\bar J\) is the current area of the cell. At atomic equilibrium,
\(\partial E_{\mathrm{el}}/\partial\w=0\), the value of \(\Pb\) does not depend on which variables
are held fixed while \(\Fb\) is varied (Hill--Mandel condition).

\paragraph{Cell reference area.}
The work done on the cell per unit change of \(\Fb\) is \(A_0\,\Pb:\mathrm d\Fb\), with \(A_0\) the
reference area of the \emph{periodic cell}. The triangles tile the current cell, so we compute
\begin{equation}
  A_0=\frac{1}{\det\Fb}\sum_e a_e .
  \label{eq:A0}
\end{equation}
This is in general different from \(\sum_e A_e\). For instance, a dislocation dipole inserted in a
periodic cell adds material half-planes between the two cores, so the summed material reference
area exceeds the cell area (in one of our test cells, by \(586\) against \(576\)). Using
\(\sum_eA_e\) in place of \eqref{eq:A0} produces a spurious offset between the prescribed and the
measured stress.

\subsection{Rotation gauge}

The energy is objective, \(E_{\mathrm{el}}(\w,\mathbf R\Fb)=E_{\mathrm{el}}(\w,\Fb)\) for every
rotation \(\mathbf R\). Every \(\Fb\) with positive determinant has a unique decomposition
\(\Fb=\mathbf R\,\Fb_{\triangle}\) with \(\Fb_\triangle\) upper triangular with positive diagonal
(the QR decomposition), so we fix the gauge
\begin{equation}
  \Fb=\begin{pmatrix}\bar F_{11}&\bar F_{12}\\0&\bar F_{22}\end{pmatrix},\qquad
  \bar F_{11},\bar F_{22}>0 .
  \label{eq:gauge}
\end{equation}
The cell then has three kinematic degrees of freedom \((\bar F_{11},\bar F_{12},\bar F_{22})\),
conjugate to \((\bar P_{11},\bar P_{12},\bar P_{22})\). The fourth component \(\bar P_{21}\) is not
an independent unknown: objectivity implies the rotational balance
\(\Pb\,\Fb^{\mathsf T}=(\Pb\,\Fb^{\mathsf T})^{\mathsf T}\), which in the gauge \eqref{eq:gauge}
reads
\begin{equation}
  \bar P_{21}\,\bar F_{11}=\bar P_{11}\,\bar F_{12}+\bar P_{12}\,\bar F_{22}.
  \label{eq:P21}
\end{equation}

\begin{proposition}\label{prop:full}
Let \(\Fb\) satisfy \eqref{eq:gauge} and let \(\Pb\) and \(\Pt\) both satisfy the rotational
balance, \(\Pb\Fb^{\mathsf T}\) and \(\Pt\Fb^{\mathsf T}\) symmetric. If
\(\bar P_{ij}=P^{\mathrm t}_{ij}\) for \((ij)\in\{11,12,22\}\), then \(\Pb=\Pt\).
\end{proposition}
\begin{proof}
\(\mathbf D=\Pb-\Pt\) has the single non-zero entry \(D_{21}=d\), and
\(\mathbf D\Fb^{\mathsf T}\) has the single non-zero entry \((\mathbf D\Fb^{\mathsf T})_{21}=d\,\bar
F_{11}\). It is symmetric only if \(d\bar F_{11}=0\), i.e.\ \(d=0\) since \(\bar F_{11}>0\).
\end{proof}

In particular, if the target is a Cauchy stress, \(\Pt=\bar J\sigt\Fb^{-\mathsf T}\), then
\(\Pt\Fb^{\mathsf T}=\bar J\sigt\) is symmetric, and matching the three free nominal components
gives the \emph{complete} Cauchy stress \(\sigb=\sigt\), including the component conjugate to the
gauged-out rotation.

%=====================================================================
\section{Nominal (Piola) stress control of a periodic cell}
\label{sec:nominal}
%=====================================================================

\subsection{Enthalpy}

Let \(\mathcal S\subseteq\{11,12,22\}\) be the set of stress-driven cell components; the others
are strain-driven (prescribed). For a prescribed nominal stress \(\Pt\) the external work is linear
in \(\Fb\), the load is dead, and the equilibria are the stationary points of the enthalpy
\begin{equation}
  H(\w,\Fb)=E_{\mathrm{el}}(\w,\Fb)-A_0\sum_{(ij)\in\mathcal S}P^{\mathrm t}_{ij}
  \big(\bar F_{ij}-\delta_{ij}\big).
  \label{eq:H}
\end{equation}
With \eqref{eq:dEdw}--\eqref{eq:macrostress}, the stationarity conditions are
\begin{align}
  \frac{\partial H}{\partial\w_a}&=\Fb^{\mathsf T}\,\frac{\partial E_{\mathrm{el}}}{\partial\x_a}=\mathbf 0
  &&\text{(atomic equilibrium)},\label{eq:eqw}\\
  \frac{\partial H}{\partial\bar F_{ij}}&=A_0\big(\bar P_{ij}-P^{\mathrm t}_{ij}\big)=0,
  \quad (ij)\in\mathcal S &&\text{(macroscopic equilibrium)}.\label{eq:eqF}
\end{align}
Since the work term is linear, \(\nabla^2H=\nabla^2E_{\mathrm{el}}\) in the joint variables
\((\w,\Fb_{\mathcal S})\): an equilibrium is stable when the elastic Hessian, with the cell
components in \(\mathcal S\) free, is positive definite. This is a weaker condition than stability
at fixed \(\Fb\) (strain control), and the difference is the source of the strain bursts discussed
in Section~\ref{sec:protocol}.

\subsection{Loading presets}

Table~\ref{tab:presets} lists the loadings implemented. A ``shear'' loading prescribes
\(P^{\mathrm t}_{12}=\alpha\) with \(P^{\mathrm t}_{11}=P^{\mathrm t}_{22}=0\), all three components
being free; uniaxial loadings prescribe one normal component. A mixed loading keeps
\(\bar F_{12}=\alpha\) prescribed and relaxes the two normal components to zero nominal stress.
An optional zero-stress relaxation of all three components at \(\alpha=0\) removes the residual
stress of the initial microstructure (e.g.\ of an inserted dipole) before loading.

\begin{table}[t]
\centering
\begin{tabular}{@{}llll@{}}
\toprule
Loading & free \(\mathcal S\) & target & conjugate strain \(\varepsilon\)\\
\midrule
nominal shear & 11, 12, 22 & \(P_{12}=\alpha\), \(P_{11}=P_{22}=0\) & \(\bar F_{12}-\bar F^{(0)}_{12}\)\\
nominal uniaxial \(y\) & 11, 12, 22 & \(P_{22}=\alpha\), others 0 & \(\bar F_{22}-\bar F^{(0)}_{22}\)\\
nominal uniaxial \(x\) & 11, 12, 22 & \(P_{11}=\alpha\), others 0 & \(\bar F_{11}-\bar F^{(0)}_{11}\)\\
Cauchy shear & 11, 12, 22 & \(\sigma_{12}=\sigma_{21}=\alpha\), \(\sigma_{11}=\sigma_{22}=0\) & \(\bar F_{12}-\bar F^{(0)}_{12}\)\\
Cauchy uniaxial \(y\) & 11, 12, 22 & \(\sigma_{22}=\alpha\), others 0 & \(\bar F_{22}-\bar F^{(0)}_{22}\)\\
Cauchy uniaxial \(x\) & 11, 12, 22 & \(\sigma_{11}=\alpha\), others 0 & \(\bar F_{11}-\bar F^{(0)}_{11}\)\\
mixed & 11, 22 & \(\bar F_{12}=\alpha\), \(P_{11}=P_{22}=0\) & ---\\
zero-stress relaxation & 11, 12, 22 & \(\Pk=\mathbf 0\) at \(\alpha=0\) & ---\\
\bottomrule
\end{tabular}
\caption{Loadings of the periodic cell. \(\Fb^{(0)}\) is the cell after the initial relaxation
step, the origin of the strain axis.}
\label{tab:presets}
\end{table}

\subsection{Nominal shear is not a pure Cauchy shear}

The nominal and Cauchy shear loadings differ by more than a change of measure. Take
\(\Pb=\tau\,\mathbf e_1\otimes\mathbf e_2+\bar P_{21}\,\mathbf e_2\otimes\mathbf e_1\) (i.e.
\(\bar P_{11}=\bar P_{22}=0\), \(\bar P_{12}=\tau\)), with \(\bar P_{21}=\tau\bar F_{22}/\bar F_{11}\)
from \eqref{eq:P21}. With \(\bar J=\bar F_{11}\bar F_{22}\), \eqref{eq:macrostress} gives
\begin{equation}
  \sigb=\frac{\tau}{\bar F_{11}\bar F_{22}}
  \begin{pmatrix}\bar F_{12}&\bar F_{22}\\ \bar F_{22}&0\end{pmatrix},\qquad
  \bar\sigma_{12}=\frac{\tau}{\bar F_{11}},\quad
  \bar\sigma_{11}=\frac{\tau\,\bar F_{12}}{\bar F_{11}\bar F_{22}} .
  \label{eq:nominalshear}
\end{equation}
Under nominal shear the cell therefore carries a normal Cauchy stress
\(\bar\sigma_{11}\approx\tau\gamma\) that grows with the shear \(\gamma=\bar F_{12}\). At the shear
strains reached before dislocation nucleation (\(\gamma\sim0.1\)--\(1\)) this is not small, and the
resolved shear stress differs from the one obtained under a pure shear stress.

%=====================================================================
\section{Cauchy (true) stress control of a periodic cell}
\label{sec:cauchy}
%=====================================================================

\subsection{The Cauchy target is a follower load}

A prescribed Cauchy stress \(\sigt\) corresponds to the deformation-dependent nominal target
\begin{equation}
  \Pt(\Fb)=\bar J\,\sigt\,\Fb^{-\mathsf T}.
  \label{eq:Ptsig}
\end{equation}
In the gauge \eqref{eq:gauge}, with
\(\Fb^{-\mathsf T}=\begin{pmatrix}1/\bar F_{11}&0\\-\bar F_{12}/(\bar F_{11}\bar F_{22})&1/\bar F_{22}\end{pmatrix}\),
the three loadings of Table~\ref{tab:presets} give the explicit targets
\begin{align}
  \sigt=\tau(\mathbf e_1\otimes\mathbf e_2+\mathbf e_2\otimes\mathbf e_1):&\quad
  P^{\mathrm t}_{11}=-\tau\bar F_{12},\quad P^{\mathrm t}_{12}=\tau\bar F_{11},\quad P^{\mathrm t}_{22}=0,
  \label{eq:Ptshear}\\
  \sigt=s\,\mathbf e_2\otimes\mathbf e_2:&\quad
  P^{\mathrm t}_{11}=0,\quad P^{\mathrm t}_{12}=0,\quad P^{\mathrm t}_{22}=s\,\bar F_{11},
  \label{eq:Ptyy}\\
  \sigt=s\,\mathbf e_1\otimes\mathbf e_1:&\quad
  P^{\mathrm t}_{11}=s\,\bar F_{22},\quad P^{\mathrm t}_{12}=0,\quad P^{\mathrm t}_{22}=0 .
  \label{eq:Ptxx}
\end{align}
A pure Cauchy shear thus requires a compressive nominal normal stress \(-\tau\bar F_{12}\) on the
cell, which cancels the normal Cauchy stress of \eqref{eq:nominalshear}. The work one-form
\(\omega=A_0\sum_{\mathcal S}P^{\mathrm t}_{ij}(\Fb)\,\mathrm d\bar F_{ij}\) is not closed:
for \eqref{eq:Ptshear}, \(\omega=A_0\tau(-\bar F_{12}\,\mathrm d\bar F_{11}+\bar F_{11}\,\mathrm d\bar F_{12})\)
and \(\mathrm d\omega=2A_0\tau\,\mathrm d\bar F_{11}\wedge\mathrm d\bar F_{12}\neq0\); likewise
\(\mathrm d\omega=A_0 s\,\mathrm d\bar F_{11}\wedge\mathrm d\bar F_{22}\) for \eqref{eq:Ptyy} and
\(-A_0 s\,\mathrm d\bar F_{11}\wedge\mathrm d\bar F_{22}\) for \eqref{eq:Ptxx}. A prescribed
deviatoric Cauchy stress therefore has no work potential, and the equilibria cannot be obtained by
minimizing a single enthalpy. (A hydrostatic target \(\sigt=-p\I\) is the exception:
\(\omega=-pA_0\,\mathrm d\bar J\) is exact.)

\subsection{Fixed-point iteration on the nominal target}

The equilibrium conditions under Cauchy control are \eqref{eq:eqw} together with
\begin{equation}
  \bar P_{ij}(\w,\Fb)=P^{\mathrm t}_{ij}(\Fb),\qquad(ij)\in\mathcal S .
  \label{eq:eqCauchy}
\end{equation}
For any frozen value \(\Pt_k\) of the target, \eqref{eq:eqw}--\eqref{eq:eqF} are the stationarity
conditions of the dead-load enthalpy \eqref{eq:H}. We therefore solve \eqref{eq:eqCauchy} by the
fixed-point iteration of Algorithm~\ref{alg:cauchy}: freeze the nominal target at the current cell,
minimize the enthalpy, re-evaluate the target at the relaxed cell, and repeat until the target no
longer changes. Each pass is warm-started from the previous minimizer, so after the first pass the
remaining passes are short.

\begin{algorithm}[t]
\caption{Cauchy stress control of a periodic cell at load level \(\alpha\)}
\label{alg:cauchy}
\begin{algorithmic}[1]
\Require configuration \((\w,\Fb)\) from the predictor; Cauchy target \(\sigt(\alpha)\)
\State \(\mathrm{tol}\gets10^{-11}+10^{-9}\max_{ij}|\sigma^{\mathrm t}_{ij}|\)
\For{\(k=1,\dots,60\)}
  \State \(\Pt_k\gets\det\Fb\;\sigt\,\Fb^{-\mathsf T}\)
  \State \((\w,\Fb)\gets\arg\min H(\,\cdot\,;\Pt_k)\) \Comment{L-BFGS from the current state, Eq.~\eqref{eq:H}}
  \State \(\Pt_{k+1}\gets\det\Fb\;\sigt\,\Fb^{-\mathsf T}\)
  \If{\(\max_{ij}|P^{\mathrm t}_{k+1,ij}-P^{\mathrm t}_{k,ij}|\le\mathrm{tol}\)} \State \Return converged \EndIf
  \If{a stress-driven \(\bar F_{ij}\) moved by more than \(\Delta F_{\max}\) in this load step}
     \State \Return limit load \Comment{Section~\ref{sec:limit}}
  \EndIf
\EndFor
\end{algorithmic}
\end{algorithm}

\paragraph{Convergence.}
Linearize about the solution \(\Fb^\ast\). Let \(\mathbb K\) be the relaxed macroscopic tangent of
the cell, i.e.\ \(\mathrm d\Pb=\mathbb K:\mathrm d\Fb\) with the atoms at equilibrium, restricted to
the components in \(\mathcal S\). A perturbation \(\delta\Fb_k\) of the cell changes the frozen target
by \(\delta\Pt_k=\mathbb L:\delta\Fb_k\) with \(\mathbb L=\partial\Pt/\partial\Fb\), and the next
minimization returns \(\delta\Fb_{k+1}=\mathbb K^{-1}\mathbb L\,\delta\Fb_k\). From
\eqref{eq:Ptshear}--\eqref{eq:Ptxx}, \(\|\mathbb L\|=O(|\sigt|)\), so the iteration contracts with
ratio
\begin{equation}
  \rho=\|\mathbb K^{-1}\mathbb L\|\sim\frac{|\sigt|}{\mu_{\mathrm{eff}}},
  \label{eq:rho}
\end{equation}
where \(\mu_{\mathrm{eff}}\) is the smallest relaxed modulus of the cell. For stresses well below
the modulus the target typically settles to the tolerance in about six passes. As the load
approaches the limit load, \(\mathbb K\) becomes singular, \(\rho\) grows without bound and the
iteration stops converging; this is detected by the limit-load guard rather than by the pass limit.

\paragraph{Exactness and the measured stress.}
At convergence, \eqref{eq:eqF} holds with \(\Pt=\bar J\sigt\Fb^{-\mathsf T}\) evaluated at the final
cell, and Proposition~\ref{prop:full} gives \(\sigb=\sigt\) for all four components. In practice
the cell Cauchy stress computed from the virial, \(\sigb=\mathbf V/(A_0\bar J)\), matches the target
to \(\sim10^{-8}\) in the controlled component, and the nominally zero components are of the same
order.

\begin{remark}
Each pass minimizes the frozen-load enthalpy, so its stability test uses
\(\nabla^2E_{\mathrm{el}}\) and omits the load stiffness \(-A_0\,\partial\Pt/\partial\Fb\) of the
follower load. That term is of relative order \(|\sigt|/\mu_{\mathrm{eff}}\), the same small
parameter as in \eqref{eq:rho}.
\end{remark}

%=====================================================================
\section{Numerical treatment of the cell variables}
\label{sec:numerics}
%=====================================================================

The unknowns of a stress-controlled relaxation are \(\mathbf z=(\w,\ \sqrt{A_0}\,\bar F_{ij})\)
for \((ij)\in\mathcal S\). Scaling the cell components by \(\sqrt{A_0}\) gives every unknown the
dimension of a length, so that the gradient components
\begin{equation}
  \frac{\partial H}{\partial(\sqrt{A_0}\,\bar F_{ij})}
  =\frac{1}{\sqrt{A_0}}\Big(\big[\mathbf V\Fb^{-\mathsf T}\big]_{ij}-A_0P^{\mathrm t}_{ij}\Big)
  =\sqrt{A_0}\,\big(\bar P_{ij}-P^{\mathrm t}_{ij}\big)
  \label{eq:gradz}
\end{equation}
have the same scaling as the atomic forces, and the limited-memory quasi-Newton (L-BFGS) Hessian
estimate is well conditioned across the two blocks. Positions are recovered as
\(\x_a=\Fb\w_a\), and the periodic image translations follow \(\Fb\) in every function evaluation.

Lagrange reduction makes the homogeneous energy bounded and periodic along lattice-invariant shears
of the cell, while the work term in \eqref{eq:H} grows linearly in \(\bar F_{12}\). For any
\(\tau\neq0\), \(H\) is therefore unbounded below: all equilibria under stress control are local
minima, separated from runaway shearing of the cell by the barrier of the energy landscape. To keep
the line search from extrapolating across that barrier, the step length of the minimizer is
capped (to \(0.5\) in units of the lattice parameter in \(\mathbf z\)).

%=====================================================================
\section{Load stepping, avalanches and the limit load}
\label{sec:protocol}
%=====================================================================

\subsection{Schedule and predictor}

The load parameter \(\alpha\) runs through a schedule \(\alpha_0,\alpha_1,\dots\); under stress
control \(\alpha\) is the prescribed stress component (\(P_{12}\), \(P_{22}\), \(P_{11}\),
\(\sigma_{12}\), \(\sigma_{22}\) or \(\sigma_{11}\)). At each step a predictor is applied before
relaxation. Strain-driven cell components follow the schedule from the relaxed initial cell
\(\Fb^{(0)}\). Stress-driven components are extrapolated from the last two converged cells,
\begin{equation}
  \bar F^{\mathrm{pred}}_{ij}=\bar F_{ij}^{(n)}+\frac{\Delta\alpha_{n+1}}{\Delta\alpha_n}
  \big(\bar F_{ij}^{(n)}-\bar F_{ij}^{(n-1)}\big),
\end{equation}
but only when step \(n\) was smooth (no reconnection and no strain burst); otherwise the predictor
is the last converged cell. The relaxation then follows (Sections~\ref{sec:nominal}--\ref{sec:cauchy}),
followed by the cascade of Delaunay reconnections of the mesh, each accepted only if it lowers the
energy functional being minimized, i.e.\ \(H\) under stress control.

\subsection{Initial relaxation at a prescribed shear strain}

A perfect periodic crystal under stress control loads elastically until the limit load, where the
cell shears away without an intermediate defected state. To start the stress-controlled run from a
dislocated microstructure, the first step can be performed under \emph{strain} control at a
non-zero shear \(\bar F_{12}=\gamma_0\), with positional noise, together with the reconnection
cascade. Beyond the homogeneous instability of the lattice (for the potentials studied here,
\(\gamma_0\gtrsim0.14\)) and in a sufficiently large cell, this relaxation nucleates a population of
dislocations. The relaxed shear stress \(\sigma_0\) of that state is then measured, and the stress
schedule is \(\alpha_k=\sigma_0+k\,\Delta\alpha\), \(k\geq1\), so that the first stress step
continues from the stress the defected crystal already carries, with no jump in load.

\subsection{Strain bursts}

Under strain control an avalanche appears as a drop of stress at fixed strain. Under force control
the stress is fixed and an avalanche appears instead as a burst of the work-conjugate strain
\(\varepsilon\) (Table~\ref{tab:presets}; for boundary tractions, Eq.~\eqref{eq:epsT} below). Let
\(\Delta\varepsilon_{\mathrm{el}}/\Delta\alpha_{\mathrm{el}}\) be the secant compliance of the last
step that showed no burst. Step \(n\) is flagged as an avalanche when
\begin{equation}
  \frac{|\Delta\varepsilon_n|-|\Delta\varepsilon^{\mathrm{el}}_n|}{|\Delta\varepsilon^{\mathrm{el}}_n|}
  \ \ge\ \theta,
  \qquad
  \Delta\varepsilon^{\mathrm{el}}_n=\Delta\varepsilon_{\mathrm{el}}\,
  \frac{\Delta\alpha_n}{\Delta\alpha_{\mathrm{el}}},
  \label{eq:burst}
\end{equation}
and, when remeshing is enabled, the reconnection cascade of the step was accepted and lowered the
enthalpy below its value after the continuous relaxation. The left-hand side of \eqref{eq:burst}
is the fractional strain burst, the force-control analogue of the fractional stress drop.

\paragraph{Energy reported.}
Under force control the functional minimized is \(H=E_{\mathrm{el}}-W_{\mathrm{ext}}\), which
decreases monotonically along a loading path dominated by the work term. The elastic energy
\(E_{\mathrm{el}}\) is reported separately: it drops when dislocations annihilate or leave the
cell and builds up as stress is stored.

\subsection{Limit load}
\label{sec:limit}

Past the limit load (the yield or Peierls stress of the current microstructure) the dead load has
no equilibrium, and the relaxation produces a runaway configuration. The step is declared a
runaway when, within one load step, a stress-driven cell component changes by more than
\(\Delta F_{\max}=0.5\) (periodic cells), or any atom moves by more than \(0.35\max(L_x,L_y)\)
(finite specimens, Section~\ref{sec:traction}). The guard is checked after the continuous
relaxation, in which case the reconnection cascade is skipped, and again after the cascade. When it
triggers, the configuration, mesh and cell at the start of the step are restored and written as
the final result, and the limit load is reported as lying in \((\alpha_{n-1},\alpha_n]\).

%=====================================================================
\section{Finite specimens: boundary tractions and force grips}
\label{sec:traction}
%=====================================================================

\subsection{Dead-load tractions}

For a finite specimen with boundary \(\partial\Omega_0\), let \(\Gamma_t\subseteq\partial\Omega_0\)
carry the nominal traction \(\mathbf t(\alpha)=\alpha\,\mathbf t_0\), measured per unit reference
length and with fixed direction (dead load). The boundary edges of the reference mesh are the edges
belonging to exactly one triangle; each carries an outward unit normal \(\mathbf n\), oriented away
from the opposite vertex of its triangle. For linear elements and a traction constant on the edge,
the consistent nodal forces are half the edge resultant at each end, so
\begin{equation}
  \f_{0,a}=\sum_{\text{edges } (ab)\subset\Gamma_t}\frac{L_{ab}}{2}\,\mathbf t_0(\text{edge}),
  \label{eq:f0}
\end{equation}
with \(L_{ab}\) the reference edge length. Four traction distributions are available:
\begin{itemize}
  \item Cartesian: \(\mathbf t_0=(t_x,t_y)\);
  \item radial outward or inward: \(\mathbf t_0=\pm\mathbf n\) on outer contours and
        \(\mp\mathbf n\) on hole faces, i.e.\ always away from (toward) the pivot of the contour;
  \item circumferential: \(\mathbf t_0=\mathbf e_z\times\mathbf r/|\mathbf r|\), with
        \(\mathbf r\) the vector from the pivot of the contour (the specimen centre, or the centre
        of the hole) to the edge midpoint (torsion).
\end{itemize}
When the loaded region is an edge strip (top, bottom, left, right), only the boundary edges whose
normal points in the direction of that face (\(\mathbf n\cdot\mathbf d\geq\tfrac12\)) are loaded;
the short side edges that a thick strip captures near the corners stay traction-free. Since
\(\mathbf n\) is the reference normal, the radial loading is a dead load, not a follower pressure.

The potential energy is
\begin{equation}
  \Pi(\x;\alpha)=E_{\mathrm{el}}(\x)-\alpha\sum_a\f_{0,a}\cdot(\x_a-\X_a),
  \label{eq:Pi}
\end{equation}
summed over the atoms with free degrees of freedom, so that equilibrium of each free atom reads
\begin{equation}
  \frac{\partial E_{\mathrm{el}}}{\partial\x_a}=\alpha\,\f_{0,a}.
\end{equation}
On atoms with one constrained component (rollers) the constrained component of the gradient is
masked after the external force has been added, so that the reaction is carried by the support.
As for the periodic cell, \(\nabla^2\Pi=\nabla^2E_{\mathrm{el}}\) and the stability under dead load is
that of the elastic Hessian on the free degrees of freedom.

\paragraph{Rigid-body modes and the isostatic anchor.}
If the specimen has no kinematic boundary condition at all, \(\Pi\) is invariant under the
rigid-body motions \(\x_a\mapsto\mathbf Q\x_a+\mathbf c\) only if the load is self-equilibrated,
\begin{equation}
  \mathbf R=\sum_a\f_{0,a}=\mathbf 0,\qquad
  M=\sum_a(\X_a-\X_c)\times\f_{0,a}=0 .
\end{equation}
Otherwise \(\Pi\) is unbounded below and the problem has no solution; such a load is rejected
(with tolerances \(|\mathbf R|\le10^{-10}\sum_a|\f_{0,a}|\) and
\(|M|\le10^{-10}\,\max(L_x,L_y)\sum_a|\f_{0,a}|\)). A self-equilibrated load still leaves the three
rigid-body modes as zero-energy directions; they are removed by a minimal isostatic anchor: the
atom closest to the centroid is fixed (\(u_x=u_y=0\)), and the atom on the same lattice row that
lies farthest from it has \(u_y=0\), which removes the rotation. Because the load is balanced, the
anchor reactions vanish at equilibrium and the anchor does not perturb the solution.

\paragraph{Conjugate strain.}
The work-conjugate strain of the dead load is
\begin{equation}
  \varepsilon=\frac{1}{A_{\mathrm{s}}}\sum_a\f_{0,a}\cdot(\x_a-\X_a),
  \label{eq:epsT}
\end{equation}
where \(A_{\mathrm s}=\sum_e|T_e|\) is the reference solid area. It satisfies
\(\alpha\,\mathrm d\varepsilon=\mathrm dW_{\mathrm{ext}}/A_{\mathrm s}\), so \((\alpha,\varepsilon)\) is a
work-conjugate pair. For a rectangle \(L_x\times L_y\) loaded by \(\mathbf t_0=\mathbf e_y\) on the
top face and \(-\mathbf e_y\) on the bottom face, the external work per unit \(\alpha\) is
\(L_x(u^{\mathrm{top}}_y-u^{\mathrm{bot}}_y)\) and \(\varepsilon=(u^{\mathrm{top}}_y-u^{\mathrm{bot}}_y)/L_y\)
reduces to the engineering strain. Strain bursts are detected with \eqref{eq:burst}.

\subsection{Rigid load-controlled grips}

A rigid grip makes all atoms of a region \(\mathcal G\) move together, \(\x_a=\X_a+\U\) for
\(a\in\mathcal G\), with the common displacement \(\U\) an additional unknown, and applies the total
force \(\alpha\mathbf F_0\) with \(\mathbf F_0=L_{\mathrm{face}}\,(t_x,t_y)\), \(L_{\mathrm{face}}\)
being the reference length of the grip's boundary edges facing outward. The potential and the
equilibrium of the grip are
\begin{equation}
  \Pi(\x,\U;\alpha)=E_{\mathrm{el}}(\x,\U)-\alpha\,\mathbf F_0\cdot\U,\qquad
  \frac{\partial\Pi}{\partial\U}=\sum_{a\in\mathcal G}\frac{\partial E_{\mathrm{el}}}{\partial\x_a}
  -\alpha\mathbf F_0=\mathbf 0 ,
\end{equation}
i.e.\ the resultant of the internal forces on the gripped atoms balances the applied force,
while the distribution of the load over the grip is whatever keeps the grip rigid. A component of
\(\U\) that carries no load can either slide freely or be held at its current value. A grip only
transmits the load, so another kinematic support (clamped, prescribed, affine or roller) is
required to react it. The conjugate strain is \(\varepsilon=\mathbf F_0\cdot\U/A_{\mathrm s}\).

Compared with distributed dead-load tractions, the grip imposes a uniform displacement of the
loaded face, as in a mechanical test, and is not sensitive to the local compliance of the face
near the corners.

\subsection{Predictor}

Pure force control has no affine predictor. The free atoms (and the grip displacements) are
extrapolated from the last two converged states,
\(\x^{\mathrm{pred}}=\x^{(n)}+\frac{\Delta\alpha_{n+1}}{\Delta\alpha_n}(\x^{(n)}-\x^{(n-1)})\), when
step \(n\) was smooth, with roller components extrapolated only along their free direction.

%=====================================================================
\section{Newton-type alternatives for the periodic cell}
\label{sec:newton}
%=====================================================================

The method of Sections~\ref{sec:nominal}--\ref{sec:cauchy} uses first derivatives only: the cell
components are additional unknowns of an L-BFGS minimization, and the Cauchy target is reached by a
fixed-point iteration on the nominal target. This section sets out Newton-type alternatives that use
second derivatives of the energy, and compares them with the present method. Apart from variant
A with the reference tangent, which an earlier stress-controlled example of the code implemented,
these are options and not part of the implementation described above.

\subsection{Second derivatives of the cell energy}

Let \(\mathbb A_e=\partial^2W/\partial\F\partial\F\) at \(\F_e\), with components
\(\mathbb A_{e,iJkL}\). From \(\F_e=\Fb\G_e\) (Eq.~\eqref{eq:FeG}) and
\(\partial\F_e/\partial w_{a,k}=\Fb\,\mathbf e_k\otimes\nabla N^e_a\), the blocks of the Hessian of
\(E_{\mathrm{el}}(\w,\Fb)\) are
\begin{align}
  K^{FF}_{ij,kl}&=\frac{\partial^2E_{\mathrm{el}}}{\partial\bar F_{ij}\,\partial\bar F_{kl}}
   =\sum_e A_e\,\mathbb A_{e,iJkL}\,G_{e,jJ}\,G_{e,lL},
  \label{eq:KFF}\\
  K^{Fw}_{ij,ak}&=\frac{\partial^2E_{\mathrm{el}}}{\partial\bar F_{ij}\,\partial w_{a,k}}
   =\sum_{e\ni a}A_e\Big(\mathbb A_{e,iJmL}\,\bar F_{mk}\,\nabla N^e_{a,L}\,G_{e,jJ}
     +P_{e,iJ}\,\delta_{jk}\,\nabla N^e_{a,J}\Big),
  \label{eq:KFw}\\
  \mathbf K^{ww}_{ab}&=\frac{\partial^2E_{\mathrm{el}}}{\partial\w_a\,\partial\w_b}
   =\Fb^{\mathsf T}\,\mathbf K^{xx}_{ab}\,\Fb ,
  \label{eq:Kww}
\end{align}
where \(\mathbf K^{xx}\) is the usual assembled stiffness with respect to positions. The second
term of \eqref{eq:KFw} is geometric: it comes from the dependence of \(\G_e\) on \(\w\). The
elastic Hessian with respect to \(\w\) and the free components of \(\Fb\) is the block matrix
\begin{equation}
  \mathbf K=\begin{pmatrix}\mathbf K^{ww}&\mathbf K^{wF}\\ \mathbf K^{Fw}&\mathbf K^{FF}\end{pmatrix},
  \qquad \mathbf K^{wF}=(\mathbf K^{Fw})^{\mathsf T}.
  \label{eq:Kfull}
\end{equation}

\paragraph{Relaxed and Cauchy--Born tangents.}
Eliminating the atoms at equilibrium gives the relaxed (effective) tangent of the cell, the Schur
complement
\begin{equation}
  A_0\,\mathbb K_{\mathrm{eff}}=\mathbf K^{FF}-\mathbf K^{Fw}\,(\mathbf K^{ww})^{-1}\,\mathbf K^{wF},
  \qquad \mathrm d\Pb=\mathbb K_{\mathrm{eff}}:\mathrm d\Fb\ \ \text{at}\ \
  \partial E_{\mathrm{el}}/\partial\w=\mathbf 0 .
  \label{eq:Keff}
\end{equation}
\(\mathbf K^{FF}/A_0\) alone is the affine (unrelaxed) tangent. For a homogeneously deformed lattice all
triangles share the same \(\F_e\), \(\Pk_e\) and \(\mathbb A_e\); since a periodic fluctuation
\(\delta\w\) changes the element gradients by amounts whose area average over the cell vanishes,
\(\mathbf K^{Fw}\delta\w=\mathbf 0\), the relaxation term vanishes, and \(\mathbb K_{\mathrm{eff}}\) reduces to
the Cauchy--Born tangent \(\mathbb A(\Fb\F_0)\), the quantity computed by the acoustic-tensor
routines of the code (\(\F_0\) maps \(\mathbb Z^2\) to the lattice). In a cell with defects, \(\mathbf
K^{ww}\) is positive definite at a stable equilibrium, so the relaxation term is positive
semidefinite and
\begin{equation}
  \mathbb K_{\mathrm{eff}}\ \preceq\ \mathbf K^{FF}/A_0 :
  \label{eq:voigt}
\end{equation}
the affine tangent is always stiffer than the relaxed one, and the difference grows with the
defect density. At the limit load \(\mathbb K_{\mathrm{eff}}\) becomes singular.

The acoustic tensor \(Q_{ik}(\mathbf n)=\mathbb C_{ijkl}n_jn_l\) built from the spatial Cauchy--Born
tangent \(\mathbb C\) gives the local criterion for loss of strong ellipticity of the homogeneous
lattice, \(\min_{\mathbf n}\det\mathbf Q(\mathbf n)=0\). It approximates the onset of homogeneous
nucleation in a perfect crystal. The stability of a defected cell under stress control is the
global condition that \(\mathbf K\) of \eqref{eq:Kfull} be positive definite. Under a Cauchy target
this condition is modified by the load stiffness \(-A_0\,\mathbb L\),
\(\mathbb L=\partial\Pt/\partial\Fb\).

\subsection{Variant A: Newton iteration on the cell only (nested)}

The atoms are relaxed at fixed \(\Fb\), and the cell is updated by a Newton step on the macroscopic
residual \(\mathbf r=\Pb-\Pt(\Fb)\) (components in \(\mathcal S\)):
\begin{equation}
  \Fb_{k+1}=\Fb_k-\mathbb T^{-1}:\mathbf r_k ,
  \label{eq:boxnewton}
\end{equation}
followed by an affine map of the atoms by \(\Fb_{k+1}\Fb_k^{-1}\) and a new relaxation. The exact
Jacobian of \(\mathbf r\) is \(\mathbb J=\mathbb K_{\mathrm{eff}}-\mathbb L\). Linearizing about the
solution, the error obeys \(\mathbf e_{k+1}=(\mathbb I-\mathbb T^{-1}\mathbb J)\,\mathbf e_k\). The choice of
\(\mathbb T\) decides the convergence:
\begin{itemize}
  \item \(\mathbb T\) = Cauchy--Born tangent at the reference state (\(\C=\I\)), used with an
    under-relaxation factor of \(0.1\) and a cap of \(0.01\) on the strain correction in the earlier
    stress-controlled example of the code. The tangent does not follow the state, and the
    iteration converges only linearly, under strong damping.
  \item \(\mathbb T\) = Cauchy--Born tangent at the current cell,
    \(\mathbb A(\Fb_k\F_0)\) with the reduced metric. This is exact for a perfect crystal
    (quadratic convergence). In a defected cell \eqref{eq:voigt} makes \(\mathbb T\) too stiff: the
    iteration converges linearly with ratio
    \(1-\lambda_{\min}(\mathbb T^{-1}\mathbb J)\), which tends to \(1\) as the limit load is
    approached and \(\mathbb K_{\mathrm{eff}}\to0\).
  \item \(\mathbb T=\mathbb J\), with \(\mathbb K_{\mathrm{eff}}\) from \eqref{eq:Keff}. The Schur
    complement needs three solves with \(\mathbf K^{ww}\) (one per free component), which one sparse
    Cholesky factorization provides. Convergence is quadratic and the Cauchy load stiffness is
    included, so no separate fixed-point loop is needed.
\end{itemize}
In every case the scheme is nested: each Newton step on the cell requires a complete relaxation of
the atoms. A remeshing cascade inside the loop has no single functional to be accepted against,
since the cell is updated by a root-finding step and not by minimization.

\subsection{Variant B: full Newton on atoms and cell}

The atoms and the cell are updated together by
\begin{equation}
  \begin{pmatrix}\mathbf K^{ww}&\mathbf K^{wF}\\ \mathbf K^{Fw}&\mathbf K^{FF}-A_0\,\mathbb L\end{pmatrix}
  \begin{pmatrix}\delta\w\\ \delta\Fb\end{pmatrix}
  =-\begin{pmatrix}\partial E_{\mathrm{el}}/\partial\w\\ A_0(\Pb-\Pt)\end{pmatrix}.
  \label{eq:fullnewton}
\end{equation}
For a nominal target \(\mathbb L=0\), and \eqref{eq:fullnewton} is the Newton step for the enthalpy
\(H\); for a Cauchy target, the load stiffness makes the matrix non-symmetric, and the follower load
is treated without an outer loop. Near a stable equilibrium the convergence is quadratic. The
difficulty is the regime this model is built to study: plastic events occur exactly where \(\mathbf
K\) loses positive definiteness, and there the Newton step is not a descent direction and may
converge to a saddle point or to an unloading branch. A robust version needs globalization, for
example a trust region with a truncated conjugate-gradient inner solver, or a modified
factorization \(\mathbf K+\mu\I\) with \(\mu\) large enough to make it positive definite. With such
safeguards the method becomes a damped Newton minimizer of \(H\), and the reconnection cascade can
again be accepted on \(H\).

\subsection{Variant C: stiffness-preconditioned L-BFGS including the cell}

The present minimizer can be given Newton-like steps while keeping descent. L-BFGS with initial
inverse Hessian \(\mathbf M^{-1}\) is equivalent to plain L-BFGS in the variables
\(\mathbf M^{1/2}\mathbf z\). Taking \(\mathbf M\) as the block stiffness \eqref{eq:Kfull}, projected to
positive semidefinite element contributions, shifted (\(\mathbf M=\mathbf K_+ +s\,\I\)) and refreshed
periodically, makes the first L-BFGS step a modified Newton step, while the line search keeps every
iterate a descent step of \(H\). The code already uses this preconditioner, with a sparse
Cholesky factorization, for the atomic unknowns under strain control; the extension consists of
adding the bordering blocks \eqref{eq:KFF}--\eqref{eq:KFw}, scaled to the variable
\(\sqrt{A_0}\,\bar F_{ij}\) of Section~\ref{sec:numerics}. The Cauchy fixed-point loop is unchanged.

\subsection{Variant D: accelerating the Cauchy loop}

Independently of the inner solver, the Cauchy iteration of Algorithm~\ref{alg:cauchy} can be
written as a root-finding problem for the cell, \(\mathbf r(\Fb)=\Pb(\Fb)-\Pt(\Fb)=0\), with
\(\Pb(\Fb)\) the relaxed stress. A secant (Broyden) update of \(\mathbb J\) built from the successive
passes, or Anderson acceleration of the fixed-point map, turns the linear contraction
\eqref{eq:rho} into superlinear convergence without computing a tangent. Since every pass after the
first is warm-started and short, the gain is limited to the passes themselves.

\subsection{Comparison}

Table~\ref{tab:newton} compares the variants with the present method (first row).

\begin{table}[ht]
\centering\small
\begin{tabular}{@{}>{\raggedright\arraybackslash}p{3.0cm}>{\raggedright\arraybackslash}p{2.5cm}>{\raggedright\arraybackslash}p{3.4cm}>{\raggedright\arraybackslash}p{2.6cm}>{\raggedright\arraybackslash}p{2.2cm}@{}}
\toprule
Method & Jacobian used & Convergence & Robust at instabilities & Remesh acceptance\\
\midrule
L-BFGS on \((\w,\Fb)\) + fixed point (present) & none (gradient only) &
  superlinear (inner), linear with ratio \(|\sigt|/\mu_{\mathrm{eff}}\) (outer) &
  yes: descent with line search & on \(H\)\\
A: cell Newton, reference CB tangent & \(\mathbb A(\F_0)\) & linear, heavy damping & partly (damping) & no functional\\
A: cell Newton, current CB tangent & \(\mathbb A(\Fb\F_0)\) & quadratic (perfect crystal); linear, \(\to1\) at limit load (defects) & partly & no functional\\
A: cell Newton, Schur tangent & \(\mathbb K_{\mathrm{eff}}-\mathbb L\) & quadratic (outer); nested relaxations & no near singular \(\mathbb K_{\mathrm{eff}}\) & no functional\\
B: full Newton & \(\mathbf K\), \(-A_0\mathbb L\) & quadratic near stable equilibria & only with trust region / shift & on \(H\) if globalized\\
C: preconditioned L-BFGS & \(\mathbf K_+ +s\I\) as preconditioner & near-Newton per iteration, fewer iterations & yes: descent with line search & on \(H\)\\
D: Broyden / Anderson outer & secant estimate of \(\mathbb J\) & superlinear (outer) & as inner solver & unchanged\\
\bottomrule
\end{tabular}
\caption{Newton-type alternatives for stress control of the periodic cell. CB: Cauchy--Born.
Variants C and D can be combined with each other and with the present formulation.}
\label{tab:newton}
\end{table}

\paragraph{Cost.}
A sparse Cholesky factorization of the \(2N\times2N\) atomic stiffness of a two-dimensional mesh
costs \(O(N^{3/2})\) with nested-dissection ordering, against \(O(N)\) for an L-BFGS iteration. The
bordering by at most three cell components adds three solves with the factor. The balance between
fewer iterations and more expensive iterations depends on the size of the cell and on how often the
factorization must be refreshed. It has to be measured on representative runs, not estimated.

\paragraph{Assessment.}
The Newton variants pay off in the elastic stretches between plastic events, where the energy is
locally convex and quadratic convergence is available. They give no advantage at the plastic events
themselves, where the Hessian is indefinite and only a descent method selects the physically
relevant branch. Variant A with a Cauchy--Born tangent, current or not, is limited by
\eqref{eq:voigt} in exactly the defected states of interest. Variants C and D keep the present
formulation, including the energy acceptance of the reconnection cascade and the detection of the
limit load, and add second-order information only where it is safe to do so.

%=====================================================================
\section{Summary}
%=====================================================================

\begin{table}[ht]
\centering\small
\begin{tabular}{@{}llll@{}}
\toprule
& unknowns & functional & equilibrium\\
\midrule
PBC, nominal stress & \(\w,\ \bar F_{ij\in\mathcal S}\) & \(E_{\mathrm{el}}-A_0\Pt:(\Fb-\I)\) & \(\bar P_{ij}=P^{\mathrm t}_{ij}\)\\
PBC, Cauchy stress & \(\w,\ \bar F_{ij\in\mathcal S}\) & frozen-target \(H\), fixed point &
  \(\bar P_{ij}=[\bar J\sigt\Fb^{-\mathsf T}]_{ij}\Rightarrow\sigb=\sigt\)\\
boundary traction & free \(\x_a\) & \(E_{\mathrm{el}}-\alpha\sum_a\f_{0,a}\cdot(\x_a-\X_a)\) &
  \(\partial E_{\mathrm{el}}/\partial\x_a=\alpha\f_{0,a}\)\\
rigid grip & free \(\x_a\), \(\U\) & \(E_{\mathrm{el}}-\alpha\mathbf F_0\cdot\U\) &
  \(\sum_{\mathcal G}\partial E_{\mathrm{el}}/\partial\x_a=\alpha\mathbf F_0\)\\
\bottomrule
\end{tabular}
\end{table}

Nominal stresses and tractions defined per unit reference length are dead loads: they derive from
a work potential linear in the kinematic unknowns, and equilibria are local minima of an enthalpy
whose Hessian is the elastic one. A prescribed Cauchy stress is a follower load without a
potential; it is enforced exactly, in all four components, by iterating the dead-load problem on
the nominal target \(\bar J\sigt\Fb^{-\mathsf T}\), a contraction with ratio of order
\(|\sigt|/\mu_{\mathrm{eff}}\) that fails only at the limit load. In all cases avalanches are
identified as bursts of the work-conjugate strain at fixed load, and the run stops at the limit
load while keeping the last equilibrium state.

\end{document}
